% ======================= 结 语 =======================
\clearpage
\chapter*{结\quad 语}
\markboth{结语}{结语}
\addcontentsline{toc}{chapter}{结语}
\vspace*{-10mm}

本书的写作，始于一九六三年秋在莫斯科大学历史系资料室的一次偶然翻阅。

当时笔者的检索目标是元代驿站制度的舆图一种。因管理员误递卷宗，笔者接触到第 117 全宗的一份移交清单。该清单共十七行，其中第十一行的编号之后用铅笔注有“另置”二字。笔者曾查询其含义，答以：自五十年代起即如此标注，不详。

三十一年之后，笔者方始了解“另置”之所指。在档案编号体系中，此类标注用于标示那些不按常规序列存放的材料，位置由移交者自定，不进入总目。

本书所写的四章，即该清单十七行中的四行。

以下三点，须向读者说明。

第一，本书不是一部揭示真相的著作。真相一词所涵盖的范围过大，笔者不敢贸然使用。本书所做的，只是重新排列一批已知的事物，指出其中一种排列更能解释现有数据。若日后另有更优的排列，笔者乐于替换。

第二，本书的每一章都能独立成立，不依赖全书结论。读者可以只读第一章，视为关于月球基地的推想；也可以只读第二章，视为关于唐代能源政策的复原。本书不主张将四章合并理解为某种阴谋。书中所描述的，从来不是阴谋，而是一个过于宏大、过于缓慢、过于沉默的系统；该系统没有意图，只有功能，没有敌人，只有阈值。

第三，也是最为重要的一点：笔者并不认为本书的结论是正确的。

笔者只认为它们是必要的。

在一个积累了足够多异常证据的学科里，假定一切正常是一种智识上的不诚实。笔者宁愿提出一个可以被推翻的解释，也不愿提出一个不会被追究的解释。错误的解释可以被证伪，含糊的解释只会自我复制。

因此，若本书最终被证明有误，那也不过是它完成自身使命的一种方式。
\begin{flushright}
{\small А.\,В.\,巴兰尼科夫}\quad{\footnotesize\color{pencil} 一九九四年五月，于莫斯科}
\end{flushright}

% ======================= 附录一 =======================
\clearpage
\chapter*{附录一\quad 本书大事年表}
\markboth{附录一　本书大事年表}{附录一　本书大事年表}
\addcontentsline{toc}{chapter}{附录一\quad 本书大事年表}

\noindent
本表以本书所主张的分期为准，与中国通行历史年表并列，以便读者对照。凡“通行系年”与“本书系年”不同者，以方括号标出。

\vspace{2mm}
{\footnotesize
\renewcommand{\arraystretch}{1.22}
\begin{longtable}{@{}p{31mm}p{55mm}p{45mm}@{}}
\toprule
\textbf{本书系年} & \textbf{事件} & \textbf{主要依据}\\
\midrule
\endfirsthead
\toprule
\textbf{本书系年} & \textbf{事件} & \textbf{主要依据}\\
\midrule
\endhead
\bottomrule
\endfoot
距今约 21{,}300 年 & 第一次生态重启（R-1） & 贝加尔湖岩芯 6.03 м 层\\
距今约 11{,}600 年 & 第二次生态重启（R-2） & 贝加尔湖岩芯 4.71 м 层\\
距今约 5{,}100 年 & 第三次生态重启（R-3） & 格陵兰冰芯 GISP2\\
约前 20000—前 221 年 & 静海七号基地运行期 & 档案 № ФВ-1963-07\\
前 221 年 & “玉山”着陆舱着陆；秦统一 & 《史记》卷六；表 1-2\\
前 221 年 & “绝地天通”（轨道电梯拆除） & 《国语·楚语下》\\
前 221—公元 80 年 & 次生黄土沉积期 & 黄土高原剖面\\
约前 1900 年 & 二里头大型建筑基址的九柱洞群 & 发掘报告\\
前 221 年以后 & 柱洞底部的环状烧结面（本书另释为“天梯”） & 档案 № ФВ-1971-22\\
630 年 & 贞观四年，关中米斗三四钱 & 《贞观政要》卷一\\
629—645 年 & 玄奘西行（密钥交接） & 《大慈恩寺三藏法师传》卷六\\
652 年 & 大雁塔始建（冷却塔） & 《长安志》卷八\\
725 年 & 一行造水运浑天仪（调节器） & 《旧唐书》卷三十五\\
755 年 & 安禄山起兵（塌缩器失控） & 《资治通鉴》卷二一七\\
756 年 & 马嵬驿（相位闭合作业） & 《长恨歌》；P.2515\\
845 年 & 会昌灭佛（记忆消除场部署） & 《旧唐书》卷十八上\\
845—1040 年 & 技术文献空白期 & 《新唐书》卷五九\\
1405—1433 年 & 郑和七下西洋（采矿作业） & 《明史》卷三〇四\\
1410 年 & 神机营立营（地面终端） & 《明史》卷八九\\
1430 年 & 宣德五年，郑和第七次出航 & 《明史》卷三〇四\\
1431 年 & 宣德六年，罢采买、停造舶 & 《明宣宗实录》卷七七\\
1559 年 & 01 号原体出生 & 铭牌 1937HL:008\\
1583 年 & “遗甲十三副”起兵 & 《清太祖实录》卷一\\
1616 年 & 后金建立（天命元年） & 《清太祖实录》卷五\\
1618 年 & “七大恨”告天，起兵略地 & 《清太祖实录》卷五；《明神宗实录》卷五六八\\
1619 年 & 萨尔浒之战 & 《明神宗实录》卷五七九\\
1626 年 & 宁远之战；努尔哈赤死 & 《明史》卷二五九\\
1644 年 & 明清交替（协议调用） & 《清世祖实录》卷一\\
\end{longtable}
}

% ======================= 附录二 =======================
\clearpage
\chapter*{附录二\quad 出土器物与档案编号索引}
\markboth{附录二　出土器物与档案编号索引}{附录二　出土器物与档案编号索引}
\addcontentsline{toc}{chapter}{附录二\quad 出土器物与档案编号索引}

\noindent
本索引按编号的字母与数字顺序排列，收录本书正文与图版所引者。凡标“未见”者，指该器物在著录后失去下落；凡标“限制”者，指其现藏机构不予公开。正文所引档案分属五套编号：苏联科学院全宗（Ф.）、VFS 基金会卷宗（ФВ-）、内府杂档（内府·）、博物馆藏品号（数字起首者）与机构报告号（МАРХИ、МГУ）。五者互不相属，故本表以编号首字为序，同一首字内再按数字排列。

须说明者，编号以 ФВ 起首的一批卷宗，并非苏联时期的全宗编号，而是一九九二年以后由 VFS 基金会（Фонд «ВФС»）异常文献部接管时所加。两套编号规则并行，故同一批材料在本索引中可能以两种形式出现。

\vspace{2mm}
{\footnotesize
\renewcommand{\arraystretch}{1.0}
\begin{longtable}{@{}p{27mm}p{40mm}p{36mm}p{24mm}@{}}
\toprule
\textbf{编号} & \textbf{器物／档案} & \textbf{现藏或出处} & \textbf{状态}\\
\midrule
\endfirsthead
\toprule
\textbf{编号} & \textbf{器物／档案} & \textbf{现藏或出处} & \textbf{状态}\\
\midrule
\endhead
\bottomrule
\endfoot
1937HL:008 & “〇一”铜铭牌 & 辽宁省博物馆 & 展出\\
1958-И-12 & 伊塞克湖木箱 & 圣彼得堡人类学与民族学博物馆 & 限制\\
71YLT14:6 & 朱书“天梯”陶片 & 中国社会科学院考古研究所 & 限制\\
87ZL:0412 & “十七刻”铜牌 & 内蒙古博物院 & 展出\\
П.2515 & 敦煌写本（“白练”条） & 法国国家图书馆 & 展出\\
Э-14 & 鄂尔多斯合金片 & 俄罗斯科学院档案馆 & 未见\\
内府·万历十七年·杂·0327 & “原体”工程残页 & 中国第一历史档案馆 & 限制\\
ФВ-1958-41 & 伊塞克湖打捞报告 & VFS 基金会异常文献部 & 限制\\
ФВ-1963-07 & 远东委员会样品分析卷宗 & VFS 基金会异常文献部 & 限制\\
ФВ-1966-09 & 赫图阿拉砖室墓 & VFS 基金会异常文献部 & 限制\\
ФВ-1971-22 & 二里头柱洞群 & VFS 基金会异常文献部 & 限制\\
ФВ-1983-15 & 龙江船厂第六号船坞 & VFS 基金会异常文献部 & 限制\\
Ф.117, оп.4, д.2187 & 静海七号图纸及退役清单 & 俄科院档案馆 & 限制\\
МАРХИ-2012-07 & 大雁塔风洞模拟 & 莫斯科建筑学院 & 未见\\
МГУ-ГЕО-2022-014 & 铼锇同位素分析 & 莫斯科大学地质系 & 限制\\
\bottomrule
\end{longtable}
}

% ======================= 附录三 =======================
\clearpage
\chapter*{附录三\quad 术语对照表（中—俄—英）}
\markboth{附录三　术语对照表}{附录三　术语对照表}
\addcontentsline{toc}{chapter}{附录三\quad 术语对照表（中—俄—英）}

\noindent
本表所收，为本书自创或重新定义的术语，共十六则。凡与通行译法不同者，以星号标出，并在“说明”栏注明取舍理由。

\vspace{2mm}
{\footnotesize
\renewcommand{\arraystretch}{1.10}
\begin{longtable}{@{}p{23mm}p{34mm}p{36mm}p{34mm}@{}}
\toprule
\textbf{中译} & \textbf{俄文} & \textbf{英文} & \textbf{说明}\\
\midrule
\endfirsthead
\toprule
\textbf{中译} & \textbf{俄文} & \textbf{英文} & \textbf{说明}\\
\midrule
\endhead
\bottomrule
\endfoot
月之民 & лунный народ & lunar folk & 前 221 年撤离月面者，自称\\
华夏神族 & божественный род Китая & the Divine Lineage & 后世对月之民及其后裔的称呼\\
生态重启 & экологический перезапуск & ecological reboot & 不译“复兴”，见译者序\\
注意力坍缩 & коллапс внимания & attention collapse & 作者自造\\
记忆消除场 & поле забвения & field of oblivion & 不译“遗忘场”\\
原体 & примарх & primarch & 作者自注借自二十世纪桌面战棋，未言其名\\
基因狂暴 & генная ярость & gene-fury & 档案原文作“发热”\\
相位脉冲弹 & фазовый импульс & phase-pulse shell & —\\
动力铠甲 & силовая броня & power armour & 明档作“甲副”\\
约束环 & кольцо удержания & confinement ring & —\\
生物相容器 & биологический фазовый сосуд & bio-phase vessel & 作者拒绝对此展开\\
重构史学 & реконструктивная история & reconstructive history & —\\
四重证据法 & метод четырёх свидетельств & four-evidence method & —\\
生态维持系统 & система поддержания экоса & ecosphere maintenance system & 全书未命名\\
矿脉 & рудное тело & ore body & 明档作“宝”\\
倒计时 & обратный отсчёт & countdown & 书末专用\\
\bottomrule
\end{longtable}
}

% ======================= 参考文献 =======================
\clearpage
\chapter*{参考文献}
\markboth{参考文献}{参考文献}
\addcontentsline{toc}{chapter}{参考文献}

\begingroup
\footnotesize
\linespread{1.25}\selectfont
\setlength{\parindent}{0pt}
\setlength{\parskip}{1.6pt}

\noindent\textbf{一、中文文献}\par
\vspace{2mm}
\begin{list}{}{\leftmargin=1.8em \itemindent=-1.8em \topsep=0pt \parsep=1.4pt \itemsep=0pt}
\item 《国语》卷十八《楚语下》，上海古籍出版社 1978 年点校本。
\item 《史记》，中华书局 1959 年点校本。
\item 《旧唐书》，中华书局 1975 年点校本。
\item 《新唐书》，中华书局 1975 年点校本。
\item 《资治通鉴》，中华书局 1956 年点校本。
\item 《明史》，中华书局 1974 年点校本。
\item 《明实录》，台湾“中研院”历史语言研究所校印本，1962 年。
\item 《清太祖高皇帝实录》，中华书局 1986 年影印本。
\item 《满文老档》，中华书局 1990 年译本。
\item 《大慈恩寺三藏法师传》，中华书局 1983 年点校本。
\item 慧立、彦悰《大唐大慈恩寺三藏法师传》，金陵刻经处本。
\item 玄奘、辩机《大唐西域记校注》，季羡林等校注，中华书局 1985 年版。
\item 马欢《瀛涯胜览》，冯承钧校注本，中华书局 1955 年版。
\item 费信《星槎胜览》，冯承钧校注本，中华书局 1954 年版。
\item 茅元仪《武备志》，明天启刻本。
\item 赵士桢《神器谱》，明万历刻本。
\item 宋应星《天工开物》，崇祯十年涂绍煃刻本。
\item 安禄山《甲胄书愤》，清嘉庆间董诰等编《全唐文》卷九六二。
\item 白居易《白氏长庆集》，文学古籍刊行社 1955 年影印本。
\item 圆仁《入唐求法巡礼行记》，上海古籍出版社 1986 年点校本。
\item 吕澂《新编汉文大藏经目录》，齐鲁书社 1980 年版。
\item 张建国《伪史学的技术伦理》，《史学理论研究》2022 年第 3 期。
\item 李文亚《郑和船队工匠配置的异常》，《海交史研究》2013 年第 2 期。
\item 《河南偃师二里头遗址三区发掘简报》，《考古》1974 年第 5 期。
\item 《陕西唐长安城郊农耕层土壤同位素分析》，《考古学报》1986 年第 2 期。
\end{list}

\vspace{5mm}
\noindent\textbf{二、俄文文献}\par
\vspace{2mm}
\begin{list}{}{\leftmargin=1.8em \itemindent=-1.8em \topsep=0pt \parsep=1.4pt \itemsep=0pt}
\item Баранников А. В. Реконструктивная история: введение в метод. М.: Изд-во МГУ, 2015.
\item Баранников А. В. Лунная прародина: опыт реконструкции. М.: Изд-во МГУ, 2018.
\item Баранников А. В. Ключевая археология Шёлкового пути. СПб.: Академическое изд-во, 2019.
\item Баранников А. В. Восстановление антиматериальной сети Чанъани. М.: Изд-во МГУ, 2019.
\item Баранников А. В. Минские бойцы: структура и характеристики. М.: Изд-во МГУ, 2020.
\item Баранников А. В. Атлас реконструкций лунных объектов. СПб., 2016.
\item Баранников А. В. «Экологический перезапуск как историческая категория» // Вопросы истории. 1973. № 9.
\item Баранников А. В. «Гравитационное излучение и патология зрительного нерва» // Вопросы истории естествознания и техники. 2004. № 3.
\item Соколов Б. П. Осадочные керны Байкала. М.: Наука, 1961.
\item Соколов В. П. Металлургия лунных базальтов. М.: Наука, 1978.
\item Кузнецов А. И. Отчёт Восточно-Сибирской экспедиции 1966 г. Новосибирск, 1967. (Архив СО РАН, ф. 12, оп. 3, д. 44.)
\item Селезнёва Н. П. «К вопросу о датировке курганов Хэтуала» // Археология Сибири. 1981. № 4.
\end{list}

\vspace{5mm}
\noindent\textbf{三、西文文献}\par
\vspace{2mm}
\begin{list}{}{\leftmargin=1.8em \itemindent=-1.8em \topsep=0pt \parsep=1.4pt \itemsep=0pt}
\item Polo, Marco. \textit{The Travels}. Trans. H. Yule, rev. H. Cordier. London, 1903.
\item Needham, Joseph. \textit{Science and Civilisation in China}. Vol. 5, Pt. 7. Cambridge, 1986.
\item Dreyer, Edward L. \textit{Zheng He: China and the Oceans in the Early Ming}. New York, 2007.
\item European Space Agency. \textit{Space Environment Report 2019}. Noordwijk, 2019.
\item USGS. \textit{Geologic Map of the Moon}, I-1234. Reston, 1979.
\end{list}
\endgroup

% ======================= 索 引 =======================
\clearpage
\chapter*{索\quad 引}
\markboth{索引}{索引}
\addcontentsline{toc}{chapter}{索引}

\noindent
（页码为本书正文页码；条目后括号内为俄文原词。）

\vspace{4mm}
{\footnotesize
\linespread{1.30}\selectfont
\setlength{\columnsep}{16pt}
\begin{multicols}{2}
\noindent
\hangindent=1.4em 安禄山（Ань Лушань）\dotfill 17、19、20、21

\hangindent=1.4em 艾瑞巴斯（Эребус）\dotfill 20

\hangindent=1.4em 白练（белая лента）\dotfill 18

\hangindent=1.4em 贝加尔湖（Байкал）\dotfill 2

\hangindent=1.4em 次生黄土（вторичный лёсс）\dotfill 5、7

\hangindent=1.4em 大雁塔（Пагода）\dotfill 13、14、15、16

\hangindent=1.4em 丹田（даньтянь）\dotfill 22

\hangindent=1.4em 帝皇（Император）\dotfill 19、20

\hangindent=1.4em 地核冷却（остывание ядра）\dotfill 5

\hangindent=1.4em 地宫储罐（подземный резервуар）\dotfill 14

\hangindent=1.4em 飞鱼服（одеяние рыбы）\dotfill 26、27

\hangindent=1.4em 广寒宫（Гуанханьгун）\dotfill 3、4、9

\hangindent=1.4em 含嘉仓（Ханьцзяцзан）\dotfill 16

\hangindent=1.4em 荷鲁斯（Гор）\dotfill 19

\hangindent=1.4em 赫图阿拉（Хэтуало）\dotfill 32、34、35、36、38

\hangindent=1.4em 记忆消除场（поле забвения）\dotfill 21、22

\hangindent=1.4em 《甲胄书愤》（«Панцирная скорбь»）\dotfill 20、21

\hangindent=1.4em 金丹（золотая пилюля）\dotfill 3、12、22

\hangindent=1.4em 绝地天通（разрыв неба）\dotfill 7、8、9

\hangindent=1.4em 昆仑山号（«Куньлунь»）\dotfill 5

\hangindent=1.4em 洛迦（Лоргар）\dotfill 20

\hangindent=1.4em 马嵬驿（Мавэй）\dotfill 17

\hangindent=1.4em 龙江船厂（верфь Лунцзян）\dotfill 29、30

\hangindent=1.4em 努尔哈赤（Нурхаци）\dotfill 32、33、34、36、37、38

\hangindent=1.4em 屏蔽环（экранное кольцо）\dotfill 4

\hangindent=1.4em 萨尔浒（Сарху）\dotfill 36、37

\hangindent=1.4em 神机营（Шэньцзиин）\dotfill 25、26、27、33

\hangindent=1.4em 生态重启（перезапуск）\dotfill 2、3

\hangindent=1.4em 十三副（тринадцать доспехов）\dotfill 34、35

\hangindent=1.4em 四神（четыре бога）\dotfill 19、20

\hangindent=1.4em 太医院（Тайиюань）\dotfill 32、33、38

\hangindent=1.4em 天宝十四载（14-й год Тяньбао）\dotfill 17、18

\hangindent=1.4em 天山基地 TK-9（база ТК-9）\dotfill 13

\hangindent=1.4em 突厥（тюрки）\dotfill 13、19

\hangindent=1.4em 万历十七年（фрагмент 1589）\dotfill 33、34、38

\hangindent=1.4em 相位脉冲（фазовый импульс）\dotfill 36

\hangindent=1.4em 玄奘（Сюаньцзан）\dotfill 11、13、14

\hangindent=1.4em 一行（Исин）\dotfill 3、5、7、16、30、34、38、39

\hangindent=1.4em 原体（примарх）\dotfill 32、33、34、35、37、38、39

\hangindent=1.4em 《长恨歌》（«Песнь о печали»）\dotfill 11、18、23

\hangindent=1.4em 郑和（Чжэн Хэ）\dotfill 28、29、30

\hangindent=1.4em 重氢（дейтерий）\dotfill 29

\hangindent=1.4em 朱雀大街（улица Чжуцюэ）\dotfill 14、16

\hangindent=1.4em 108 坊（108 кварталов）\dotfill 14、16

\hangindent=1.4em 四重证据法（четыре свидетельства）\dotfill 序 vi

\hangindent=1.4em 倒计时（обратный отсчёт）\dotfill 全书

\end{multicols}
}

% ======================= 译 后 记 =======================
\clearpage
\chapter*{译\quad 后\quad 记}
\markboth{译后记}{译后记}
\addcontentsline{toc}{chapter}{译后记}

\noindent
本书的翻译历时十四个月，其间三易其稿。

第一稿完成于二〇二二年春，由译者自行废弃。其原因在于，该稿在处理方法上侧重于替作者辩护：凡读者可能质疑之处，均加“译者按”一句，说明此说未必成立。此稿曾送请一位前辈审阅，其意见是：这样译等于提前替读者把书读完了。

第二稿删去全部辩解，仅保留必要的术语说明，结果亦被废弃。其原因在于，不设辩解，則近于为作者背书，同样不妥。

第三稿即现在所见的本子。其处理原则为：正文照译，不作改动；存疑之处一律入注；注释中的“译者注”只陈述事实，不发表意见。

译者以为，读者所需要的不是判断的替代，而是完整的信息，以及完整信息所必然带来的那种不适。

最后须说明一处细节。本书第三章、第四章中，“四十分钟”一语共出现四次。译至第三章时，译者曾就此作过一次统计，所得结论为四次。

其后核对俄文原版，发现第五次见于全书索引的最后一行，条目如下：

\vspace{2mm}
\begin{center}
{\small\itshape Обратный отсчёт（倒计时）\dotfill 全书}
\end{center}
\vspace{2mm}

页次一栏所填，不是页码，而是“全书”二字。

\vspace{8mm}
\begin{flushright}
{\small 李文亚\ \ 张建国}\\[2pt]
{\footnotesize\color{pencil} 二〇二三年十二月，于北京}
\end{flushright}

% ======================= 勘 误 汇 总 =======================
\clearpage
\chapter*{初版至三版勘误汇总}
\markboth{初版至三版勘误汇总}{初版至三版勘误汇总}
\addcontentsline{toc}{chapter}{初版至三版勘误汇总}

\noindent
本表所列，是初版至三版发现并在本版改正的条目。页次据第三版（二〇二四年九月），行次从略。经考订而仍有疑问的，另加按语，不作径改。

\vspace{4mm}
{\footnotesize
\renewcommand{\arraystretch}{1.22}
\begin{longtable}{@{}p{17mm}p{58mm}p{58mm}@{}}
\toprule
\textbf{页} & \textbf{误} & \textbf{正}\\
\midrule
\endfirsthead
\toprule
\textbf{页} & \textbf{误} & \textbf{正}\\
\midrule
\endhead
\bottomrule
\endfoot
ii & 丛书页“viii 唐长安城反物质管网复原”署沃兹涅先斯基，“x 郑和舰队南极补给点考”署李文亚 & 两卷均署巴兰尼科夫（与正文及参考文献一致）\\
2 & 两层夹层距今约一万一千二百年与两万四千六百年 & 距今约一万一千六百年与两万一千三百年（与表 1-1 层位对应）\\
5 & 关于该船体本身……船体至今仍在地球轨道上 & 母船在轨（“月球”），着陆舱已解体，残片散见黄土\\
6 & 抛射物总量 1.9×10\textsuperscript{12} t；厚达数十米的次生黄土 & 9.4×10\textsuperscript{12} t；压实后十余米\\
9 & 同期在籍人口仅恢复了一成左右 & 仅恢复至隋大业年间的三成四\\
9 & 他们以青铜代替合金 & 以铁和青铜代替合金（战国铁器已普及，所失者为合金工艺）\\
12 & 图 2-1 比例尺“10 м” & “5 м”（塔高 64.5 米）\\
13、14 & 两段注文脱去注符，误入正文 & 已归其位\\
17 & 其时安禄山已死四十余年 & 六十年\\
17 & 是年七月，长安上空出现一道“裂缝” & 是年六月（下引本证亦作六月）\\
17 & 至德元载的户口统计……全国平均减少 36\% & 广德二年……67\%（唐代无至德元载户口统计）\\
18 & 距事发已五十一年 & 五十年\\
22 & 《明史》卷八十九《兵志四》 & 卷八十九《兵志一》；卷九十二《兵志四》\\
22 & 一次完整的打击循环约需十七分钟 & 十七刻，合四小时零五分\\
24 & 总载荷 1.86 kN & 1.06 kN（装具 38 公斤＋乘员约 70 公斤）\\
24 & 蒸汽推力径直由汽缸直径与压力反算 & 单缸 2.67 kN，经髋部放大器四点六倍而得\\
26 & 长 44 丈 4 尺（138.6 м）、宽 18 丈（56.2 м） & 142.1 м、57.6 м（与本书所采明尺 32 厘米相合）\\
27 & 一九八三年……船坞长约 100 米 & 二〇〇三年……约 240 米（100 米装不下 142 米宝船）\\
27、28 & 宣德五年，郑和最后一次归航 & 宣德五年六月出航，宣德八年七月归国\\
31 & 铭牌背面刻“E-13 / GENE-PRIMARCH…” & “G-13 /…”（与后文缺陷编号一致）\\
32 & 表 4-1 题“遗甲十三副复原清单”，首列甲—01 至甲—13 & 题改为构件与残件对照，首列件—01 至件—13\\
34 & 朝鲜《光海君日记》卷一九六记 1626 年事 & 《仁祖实录》卷十四（光海君 1623 年已废）\\
37 & 年表：1409 年 神机营立营 & 1410 年（永乐八年）\\
37 & 年表：1431 年 宣德罢采买、停造舶 & 1430 年出航，1431 年罢舶，析为两行\\
37 & 年表：静海七号运行期止于前 3000 年 & 止于前 221 年（与撤离之年合）\\
37 & 年表：次生黄土沉积期“前 221—前 200 年” & “前 221—前 79 年”（正文作约三百年）\\
37 & 年表缺“七大恨” & 补 1618 年一行\\
41 & 清·严可均辑《全唐文》 & 清嘉庆间董诰等编《全唐文》\\
41 & 附录二“按编号首字母排序”，漏收 ФВ-1983-15 与内府杂档 & 改按字母与数字混排，补收两种编号\\
\bottomrule
\end{longtable}
}

\vspace{3mm}
{\footnotesize
\textbf{按}　第三版第三十一页“二十一条命，十三套铁”一语，经复核无误，不改。\\
\phantom{\textbf{按}　}第三版所引“十七刻”铜牌，其数合四小时零五分；旧译本径作“十七分钟”，系译者换算之误，已入上表。\\
}

\vspace{8mm}
\begin{center}
{\small\color{pencil}
——\quad 全书终\quad ——}
\end{center}
